\scanpage{8}

La commutation aux $\underleftarrow{\lim}$ adiques peut aussi s'exprimer ainsi ;

(L) Pour tt \deleted{\illegible{1}} $T$ \inserted{sur $S$} spectr. d'un anneau local noeth. \deleted{complet sur $S$}\inserted{(qu'on peut supposer \uline{complet}, à corps rés. alg. clos.)}, et tel que $Y_T$ majore un voisinage ouvert induit de la fibre spéciale de $X_T$, on a $\deleted{\unsure{Y_T|X_T}}X_T=Y_T$.

\uline{Corollaire} Cas $X$ propre sur $S$ ; \uuline{O.K.}

Si $X$ est un groupe, $Y$ un sous-groupe, on voit que $Y_T$ \deleted{contient}\inserted{majore} un voisinage ouvert induit de la section unité, donc $Y_T\to X_X$ est \inserted{un \uline{mono}} étale, \deleted{\unsure{donc}} étant une imm \deleted{fermée}, est une immersion ouvert\unsure{\deleted{e}|e} \deleted{et fermée}. \deleted{\unsure{Donc dans ce cas|Donc dans le cas}} \inserted{De plus, la vérification $Y_T=X_T$ \unsure{est la|est|se fait} } \deleted{même} chose ``ensembliste'', donc se ramène au cas où $T$ est le spectre d'un anneau de valuation discrète. Donc

\uline{Corollaire} Dans ce cas, la condition envisagée équivaut à la suivante : Pour tt $T$ sur $S$ spectr. d'un anneau \deleted{local noeth} \deleted{de \unsure{valuation}}\inserted{\uline{de valuation discrète}} (qu'on peut supposer complet à corps rés. alg clos), tel que $Y_T$ soit un sous-groupe ouvert et fermé de $X_T$ qui contient la fibre \deleted{générique}\inserted{spéciale}, on a $Y_T=X_T$. Cette condition est vérifiée dans chacun des cas suivants.

a) $X$ \uline{à fibres connexes}

b) \uline{Il existe une famille de ss-\deleted{groupes finis}\inserted{schémas}} $\deleted{\unsure{H}}G_i^{\unsure{?}}$ de $G$, finis sur $S$, et tels que pour chaque $s\in S$, la réunion ensembliste des $G_{i,\bar s}\cdot G_{\bar s}^{0}$ soit $G_{\bar s}$.
% archive sheet number: 7
